\documentclass[11pt]{article}
\usepackage[margin=2.2cm]{geometry}
\usepackage{amsmath,amssymb,booktabs,array}
\usepackage[T1]{fontenc}
\usepackage{lmodern}
\setlength{\parskip}{0.5em}\setlength{\parindent}{0pt}
\title{Toroidal-fit coherence on the OMAHA array:\\derivation and the $n$ table}
\author{MASTU\_hazard notes}
\date{}
\begin{document}
\maketitle

\section{Definition (as coded)}
For one frequency bin $f$ and one time slice $t$, let $S_c$ be the complex
spectrum of coil $c$ (\texttt{specgram(\ldots, mode='complex')}),
$c=1,\dots,N_c$, with $N_c=5$ OMAHA slow coils at toroidal angles
\[
\varphi_c = [\,36.8,\;26.5,\;4.0,\;-5.0,\;-23.0\,]^\circ .
\]
\texttt{Spectra.n\_detection} (\texttt{egio/saddle\_data.py}) forms, for each searched $n$,
\[
Z_n = \sum_{c=1}^{N_c} e^{\,i n\varphi_c}\, S_c ,
\qquad
\texttt{power} = \frac{1}{N_c}\sum_c |S_c|^2 ,
\qquad
\texttt{coherence}_n = \frac{|Z_n|^2}{N_c^2\;\texttt{power}} .
\]
Substituting \texttt{power}:
\[
\boxed{\;\mathrm{coh}_n = \frac{\bigl|\sum_c e^{\,i n\varphi_c} S_c\bigr|^2}{N_c \sum_c |S_c|^2}\;}
\]
The bin is assigned to $n^\ast = \arg\max_n |Z_n|$ over the searched list
$\{0,\pm1,\dots,\pm4\}$ (equivalently the largest $\mathrm{coh}_n$, since the
denominator does not depend on $n$). The coherence gate then keeps the bin
only if $\mathrm{coh}_{n^\ast}\ge$ \texttt{coh\_min}.
\textbf{The argmax comes before the gate.}

\section{Bounds: why a real mode scores 1}
By Cauchy--Schwarz, with $a_c = e^{in\varphi_c}$ and $|a_c|=1$,
\[
\Bigl|\sum_c a_c S_c\Bigr|^2 \le \Bigl(\sum_c |a_c|^2\Bigr)\Bigl(\sum_c |S_c|^2\Bigr) = N_c\sum_c|S_c|^2
\quad\Rightarrow\quad 0\le \mathrm{coh}_n\le 1 ,
\]
with equality iff $S_c = \lambda\,\overline{a_c} = \lambda\, e^{-in\varphi_c}$
for one complex $\lambda$ common to all coils.

A rotating mode $B_c(t)=A\cos(\omega t - m\varphi_c + \alpha)$ has, in its
positive-frequency bin, $S_c \propto \tfrac{A}{2}\,e^{i\alpha}\,e^{-im\varphi_c}$
(identical gain, area and window on every coil assumed). That is exactly the
equality case for $n=m$: \textbf{a clean mode scores $\mathrm{coh}_m = 1$ at
its own $n$.}

\section{A pure mode scored at another $n$: the array factor}
Put $S_c = \lambda e^{-im\varphi_c}$ into the boxed formula and score it at $n$:
\[
\mathrm{coh}_n
= \frac{|\lambda|^2\,\bigl|\sum_c e^{\,i(n-m)\varphi_c}\bigr|^2}{N_c\cdot N_c|\lambda|^2}
= \frac{\bigl|\sum_c e^{\,ik\varphi_c}\bigr|^2}{N_c^2},
\qquad k = n-m .
\]
It depends only on the difference $k=n-m$, not on $n$ or $m$ separately, and
$|{\textstyle\sum} e^{-ik\varphi_c}| = |{\textstyle\sum} e^{ik\varphi_c}|$, so it
is even in $k$. The common mode is the case $m=0$: a signal every coil sees
identically scores $\mathrm{coh}_n = |\sum_c e^{in\varphi_c}|^2/25$.

\subsection*{Worked numbers}
Per-coil terms $\cos k\varphi_c$, $\sin k\varphi_c$:

\begin{center}\small
\begin{tabular}{c r *{4}{r r}}
\toprule
coil & $\varphi_c$ [deg] & \multicolumn{2}{c}{$k=1$} & \multicolumn{2}{c}{$k=2$} & \multicolumn{2}{c}{$k=3$} & \multicolumn{2}{c}{$k=4$}\\
 & & cos & sin & cos & sin & cos & sin & cos & sin\\
\midrule
1 & +36.8 & +0.8007 & +0.5990 & +0.2823 & +0.9593 & -0.3486 & +0.9373 & -0.8406 & +0.5417 \\
2 & +26.5 & +0.8949 & +0.4462 & +0.6018 & +0.7986 & +0.1822 & +0.9833 & -0.2756 & +0.9613 \\
3 & +4.0 & +0.9976 & +0.0698 & +0.9903 & +0.1392 & +0.9781 & +0.2079 & +0.9613 & +0.2756 \\
4 & -5.0 & +0.9962 & -0.0872 & +0.9848 & -0.1736 & +0.9659 & -0.2588 & +0.9397 & -0.3420 \\
5 & -23.0 & +0.9205 & -0.3907 & +0.6947 & -0.7193 & +0.3584 & -0.9336 & -0.0349 & -0.9994 \\
\midrule $\Sigma$ & & +4.6099 & +0.6371 & +3.5539 & +1.0041 & +2.1361 & +0.9360 & +0.7499 & +0.4372 \\
\bottomrule
\end{tabular}
\end{center}

Then $\mathrm{coh} = (\Sigma\cos^2 + \Sigma\sin^2)/25$. For example, $k=1$:
$(4.6099^2 + 0.6371^2)/25 = (21.251 + 0.406)/25 = 21.657/25 = 0.866$.

\begin{center}
\begin{tabular}{c r r r r}
\toprule
$|k|=|n-m|$ & $\sum\cos k\varphi_c$ & $\sum\sin k\varphi_c$ & $|\sum e^{ik\varphi_c}|^2$ & $\mathrm{coh}$ \\
\midrule
0 & +5.0000 & +0.0000 & 25.000 & 1.000 \\
1 & +4.6099 & +0.6371 & 21.657 & 0.866 \\
2 & +3.5539 & +1.0041 & 13.638 & 0.546 \\
3 & +2.1361 & +0.9360 & 5.439 & 0.218 \\
4 & +0.7499 & +0.4372 & 0.753 & 0.030 \\
5 & -0.2498 & -0.3194 & 0.164 & 0.007 \\
6 & -0.6541 & -1.0574 & 1.546 & 0.062 \\
\bottomrule
\end{tabular}
\end{center}
(For $k<0$ the sine column changes sign; $\mathrm{coh}$ is unchanged.)

\textbf{Why it falls so slowly.} The five coils span only
$36.8-(-23.0)=59.8^\circ$ of the torus. Across that aperture $e^{ik\varphi}$
turns by $k\times 59.8^\circ$; it takes $k\approx 360/59.8\approx 6$ for a
full turn and a cancellation. So neighbouring $n$ are barely distinguished:
the first near-null is at $|k|=5$ (0.007), and $|k|=1$ still scores 0.866.

\section{The table by mode}
Each row: a pure mode of toroidal number $m$, scored at every searched $n$.
The diagonal (bold) is the mode's own $n$; the argmax picks it when there is
no noise.

\begin{center}\small
\begin{tabular}{c *{9}{c}}
\toprule
 & \multicolumn{9}{c}{scored at $n$}\\
mode $m$ & 0 & $+1$ & $-1$ & $+2$ & $-2$ & $+3$ & $-3$ & $+4$ & $-4$\\
\midrule
$0$ & \textbf{1.000} & 0.866 & 0.866 & 0.546 & 0.546 & 0.218 & 0.218 & 0.030 & 0.030 \\
$+1$ & 0.866 & \textbf{1.000} & 0.546 & 0.866 & 0.218 & 0.546 & 0.030 & 0.218 & 0.007 \\
$+2$ & 0.546 & 0.866 & 0.218 & \textbf{1.000} & 0.030 & 0.866 & 0.007 & 0.546 & 0.062 \\
$+3$ & 0.218 & 0.546 & 0.030 & 0.866 & 0.007 & \textbf{1.000} & 0.062 & 0.866 & 0.098 \\
\bottomrule
\end{tabular}
\end{center}
A negative $m$ (opposite rotation) is the same row with $n\to -n$.
With \texttt{absolute\_n} the $|n|$ trace pools bins assigned to $+|n|$ and $-|n|$.

\section{What follows}
\begin{enumerate}
\item \textbf{Aliasing between neighbours.} A pure $n=1$ scores 0.866 at
  $n=2$ and at $n=0$; a pure $n=2$ scores 0.866 at $n=1$ and $n=3$. The
  genuine $n$ wins the argmax by a margin of only $1-0.866=0.134$.
\item \textbf{The gate cannot remove a flipped bin.} Noise can push a bin's
  argmax to the neighbour. The coherence that bin then carries is about
  0.866, above a 0.8 gate, so it passes. Raising the gate above 0.866 rejects
  pure aliases but also starts rejecting genuine modes that noise has pulled
  below 1. This is why leakage is tested on frequency (Section~6), not by
  the gate.
\item \textbf{A genuine $n=2$ passes a 0.8 gate.} It scores 1 at $n=2$, not
  0.866. The gate only matters when noise lowers it.
\item \textbf{$n=0$ is searched as a sink, never reported.} Common-mode
  signal (baseline drift, pickup every coil sees alike) scores 1.000 at $n=0$
  but 0.866 at $n=\pm1$. Without $n=0$ in the list the argmax would send it
  to $n=\pm1$ and it would pass a 0.8 gate. With $n=0$ present the argmax
  claims it for $n=0$, and the $|n|\ge1$ masks exclude it.
\item \textbf{Noise floor.} Pure coil noise gives $E[\mathrm{coh}_n]=1/N_c=0.2$
  at every $n$, exactly. The assigned value (maximum over the nine searched
  $n$) averages about $2/N_c$. Derivation and tail probabilities in
  Section~7.
\end{enumerate}

\section{Leakage test on frequency}
Under rigid rotation a mode of number $n$ appears at $f \approx n f_{\mathrm{rot}}$.
An $n=1$ bin flipped into the $|n|=2$ trace keeps the $n=1$ frequency:
$F_a(|n|{=}2)/F_a(|n|{=}1)\approx 1$, at the same time. A genuine $n=2$, or an
overtone of $n=1$, sits at $\approx 2$. \texttt{check\_n\_leakage.py} applies
this; \texttt{run\_master.py -{}-harmonics} then separates an overtone
(phase-locked, $hp\to1$) from a separate $n=2$ ($hp\sim 1/\sqrt{N_\mathrm{eff}}$).


\section{Noise floor}
\subsection{Exact mean}
Write $a_c=e^{\,in\varphi_c}$, so $|a_c|=1$. Expand the numerator:
\[
\Bigl|\sum_c a_c S_c\Bigr|^2
= \sum_c |S_c|^2 \;+\; \sum_{c\neq d} a_c a_d^{*}\, S_c S_d^{*} .
\]
Divide by the denominator $N_c\sum_c|S_c|^2$:
\[
\mathrm{coh}_n = \frac{1}{N_c}
\;+\; \frac{1}{N_c}\,\frac{\sum_{c\neq d} a_c a_d^{*}\, S_c S_d^{*}}{\sum_e |S_e|^2} .
\]
The first term is $1/N_c$ for \emph{every} realisation. For the second,
assume each coil's noise has a uniformly random phase, independent of the
other coils. Rotating one coil, $S_c\to e^{i\theta}S_c$, leaves
$\sum_e|S_e|^2$ unchanged and multiplies every cross term containing $S_c$
by $e^{\pm i\theta}$. Averaging over $\theta$ sends each cross term to zero:
\[
E\!\left[\frac{S_c S_d^{*}}{\sum_e |S_e|^2}\right] = 0 \quad (c\neq d)
\qquad\Longrightarrow\qquad
\boxed{\;E[\mathrm{coh}_n] = \frac{1}{N_c} = 0.2\;}
\]
for every $n$. This needs neither Gaussian noise nor equal noise on the
coils; only independent, phase-random noise. A spectrogram bin away from
DC satisfies the phase condition well.

\subsection{Distribution and tail}
Assume pure noise: the $S_c$ independent, circular complex Gaussian
$\mathcal{CN}(0,\sigma^2)$, with the same $\sigma$ on every coil. Then each
noise-only bin's $\mathrm{coh}_n$, at one fixed $n$, is one draw from a Beta
distribution. Four steps:

\begin{enumerate}
\item \textbf{Unit steering vector.} With
  $\hat a = (e^{in\varphi_1},\dots,e^{in\varphi_{N_c}})/\sqrt{N_c}$, $|\hat a|=1$,
  \[
  \mathrm{coh}_n = \frac{|\hat a^{\dagger} S|^2}{|S|^2},
  \]
  the fraction of the bin's total coil power lying along the $n$ direction.
\item \textbf{Rotate.} Take a unitary $U$ whose first row is $\hat a^{\dagger}$
  and set $T=US$. White circular Gaussian noise is unchanged in distribution
  by a unitary rotation, so $T_1,\dots,T_{N_c}$ are again independent
  $\mathcal{CN}(0,\sigma^2)$. The numerator is $|\hat a^{\dagger}S|^2=|T_1|^2$;
  the denominator is $|S|^2=\sum_k|T_k|^2$, since rotation preserves length.
\item \textbf{Gammas.} Each $|T_k|^2$ is exponential, i.e.\
  $\mathrm{Gamma}(1,\sigma^2)$. So $X=|T_1|^2\sim\mathrm{Gamma}(1)$ and
  $Y=\sum_{k\ge2}|T_k|^2\sim\mathrm{Gamma}(N_c-1)$, independent of each other.
\item \textbf{Beta.} For independent Gammas of common scale,
  $X/(X+Y)\sim\mathrm{Beta}(1,N_c-1)$, and the ratio is independent of
  $X+Y$. Hence
  \[
  \mathrm{coh}_n \sim \mathrm{Beta}(1,\,N_c-1):\quad
  p(x) = (N_c-1)(1-x)^{N_c-2} = 4(1-x)^3,
  \]
  \[
  E = \frac{1}{N_c} = 0.2,\qquad
  \sigma = \sqrt{\frac{N_c-1}{N_c^2(N_c+1)}} = 0.163,\qquad
  P(\mathrm{coh}_n > x) = (1-x)^{N_c-1}.
  \]
  The independence of the ratio from $X+Y$ is also why
  $E[\text{num}]/E[\text{den}]$ happens to give the right mean in this case;
  Section~7.1 shows $1/N_c$ does not need Gaussian noise at all.
\end{enumerate}
The mean is a poor guide: the spread is nearly as large as the mean.

\begin{center}
\begin{tabular}{lcc}
\toprule
Pure noise, per bin (400k samples) & Simulated & Theory \\
\midrule
$E[\mathrm{coh}_n]$, any $n$ & 0.200 & $1/N_c = 0.2$ \\
$\sigma(\mathrm{coh}_n)$ & 0.164 & 0.163 \\
$P(\mathrm{coh}_n > 0.8)$ & 0.0017 & $0.2^4 = 0.0016$ \\
$E[\max_n \mathrm{coh}_n]$, nine $n$ & 0.40 & none derived \\
$P(\max_n \mathrm{coh}_n > 0.8)$ & 0.011 & none derived \\
\bottomrule
\end{tabular}
\end{center}
What the Beta result does \emph{not} say:
\begin{itemize}
\item \textbf{Not independent across $n$.} $\mathrm{coh}_n$ and
  $\mathrm{coh}_{n'}$ in one bin use the same five numbers and overlapping
  steering vectors (the 0.866, 0.546, \ldots\ of Section~3). Each is
  $\mathrm{Beta}(1,4)$ on its own, but they are correlated, so the maximum
  over $n$ has no simple closed form; its values above are from simulation.
\item \textbf{Not independent across bins.} Hann windows at 50\% overlap
  correlate neighbouring time slices ($\rho=1/6$), and Hann leakage
  correlates neighbouring frequency bins. Bins passing a gate on noise
  arrive in clusters; the count in Section~7.4 is an average.
\item \textbf{Not with a mode present.} A mode adds a fixed vector to $S$;
  $X$ becomes noncentral and $\mathrm{coh}_n$ is pulled toward 1 by an
  amount set by the signal-to-noise ratio. $\mathrm{Beta}(1,N_c-1)$ is the
  no-mode case only.
\item \textbf{Not with unequal coil noise.} Step 2 then fails: the rotation
  changes the noise distribution. The mean stays $1/N_c$ (Section~7.1) but
  the shape changes (Section~7.3).
\end{itemize}

\subsection{When the floor changes}
\begin{itemize}
\item \textbf{Unequal coil noise.} The mean stays $1/N_c$ (Section~7.1), but
  the tail changes. With one coil three times noisier,
  $P(\mathrm{coh}_n>0.8)$ falls to $0.0006$: the noisy coil dominates and
  no $n$ fits it.
\item \textbf{Correlated noise.} Pickup common to several coils keeps cross
  terms alive, and the floor rises at whichever $n$ matches the pickup
  pattern. Pure common mode lands at $n=0$ (Section~5, item~4).
\end{itemize}

\subsection{Consequence for detection}
A 1--50\,kHz band holds about $49\,000/390.6\approx125$ frequency bins. With
pure noise and a 0.8 gate, about $125\times0.011\approx1.4$ bins per time
slice pass at some $n$. Coherence alone does not reject noise bin by bin;
the amplitude threshold does.

\end{document}
